% =============================================================================
\section{Dieléctricos lineales, homogéneos e isótropos}\label{sec:lineales}
% =============================================================================

Pasamos ahora a los dieléctricos ordinarios, en los que la polarización no es permanente sino la respuesta
del material al campo que actúa sobre él. Para cerrar el sistema formado por \eqref{eq:gauss-D} y
$\rot\vE=\vcero$ hace falta una relación entre $\vP$ y $\vE$, que depende del material y debe obtenerse de su
estructura microscópica o de la experiencia.

\subsection{La relación constitutiva}\label{sec:constitutiva}

Desde el punto de vista macroscópico, si $\vP$ se anula con el campo y este es débil frente a los campos
atómicos ($\sim\SI{e11}{V/m}$), cada componente de la polarización puede desarrollarse en potencias de las del campo,
\begin{equation}\label{eq:constitutiva-general}
  P_i(\vr)=\sum_{j}\alpha_{ij}(\vr)\,E_j(\vr)+\sum_{j,k}\beta_{ijk}(\vr)\,E_j(\vr)E_k(\vr)+\cdots,
  \qquad i,j,k\in\{x,y,z\},
\end{equation}
donde los coeficientes son propiedades del material en el punto $\vr$. Es una relación \emph{local}: la
polarización en $\vr$ solo depende del campo en $\vr$.

\begin{definicion}[Medio lineal, homogéneo e isótropo]\label{def:LHI}
Un dieléctrico es
\begin{enumerate}
  \item \emph{lineal} si en \eqref{eq:constitutiva-general} basta el primer término; se escribe entonces
        \begin{equation}\label{eq:chi-tensor}
          P_i(\vr)=\eo\sum_j\chi_{ij}(\vr)\,E_j(\vr),
        \end{equation}
        y la matriz adimensional $(\chi_{ij})$ es el \emph{tensor de susceptibilidad eléctrica};
  \item \emph{homogéneo} si $\chi_{ij}$ no depende de $\vr$: es una propiedad intrínseca del material, la misma en
        todos sus puntos;
  \item \emph{isótropo} si no tiene direcciones privilegiadas: $\chi_{ij}=\chi\,\delta_{ij}$, es decir,
        $\chi_{xx}=\chi_{yy}=\chi_{zz}=\chi$ y $\chi_{ij}=0$ para $i\neq j$.
\end{enumerate}
En un medio \LHI{}, por tanto,
\begin{equation}\label{eq:chi}
  \boxed{\;\vP(\vr)=\eo\chi\,\vE(\vr)\;},
\end{equation}
con $\chi>0$ una constante, la \emph{susceptibilidad eléctrica} del material.
\end{definicion}

\begin{observacion}[Por qué $\chi>0$]\label{obs:chi-positiva}
Los dos mecanismos de la \cref{sec:microscopico} producen dipolos en el sentido del campo: el desplazamiento de
la nube electrónica y la orientación de los dipolos permanentes. Por eso $\vP$ tiene el sentido de $\vE$ y la
susceptibilidad es positiva. En un cristal no cúbico la respuesta depende de la dirección y $(\chi_{ij})$ no es
proporcional a la identidad, aunque sí simétrica y definida positiva~\cite{Landau}.
\end{observacion}

\subsection{Permitividad}\label{sec:permitividad}

En un medio \LHI{} el desplazamiento es proporcional al campo:
\[
  \vD=\eo\vE+\vP=\eo\vE+\eo\chi\vE=\eo(1+\chi)\,\vE .
\]

\begin{definicion}[Permitividad]\label{def:permitividad}
En un dieléctrico \LHI{} se definen la \emph{permitividad dieléctrica} $\varepsilon$ y la \emph{permitividad relativa}
(o constante dieléctrica) $\er$,
\begin{equation}\label{eq:epsilon}
  \varepsilon\defeq\eo(1+\chi),\qquad \er\defeq\frac{\varepsilon}{\eo}=1+\chi ,
\end{equation}
de modo que
\begin{equation}\label{eq:D-epsilon}
  \boxed{\;\vD=\varepsilon\,\vE,\qquad \vP=(\varepsilon-\eo)\,\vE\; .}
\end{equation}
Como $\chi>0$, $\varepsilon>\eo$ y $\er>1$. Las unidades de $\varepsilon$ son las de $\eo$, $\si{F/m}$; $\er$ es adimensional
(algunos textos la llaman $\kappa$ o $K$~\cite{Purcell}).
\end{definicion}

\begin{table}[t]
\centering
\caption{Permitividades relativas de algunos materiales a temperatura ambiente~\cite{Griffiths}.}
\label{tab:permitividades}
\begin{tabular}{@{}lc@{\qquad}lc@{}}
\toprule
Material & $\er$ & Material & $\er$\\
\midrule
Vacío & $1$ (exacto) & Diamante & $5{,}7$\\
Aire (seco) & $1{,}00054$ & Sal común (NaCl) & $5{,}9$\\
Benceno & $2{,}28$ & Silicio & $11{,}7$\\
Metanol & $33{,}0$ & Agua (\SI{20}{\celsius}) & $80{,}1$\\
\bottomrule
\end{tabular}
\end{table}

\subsection{Consecuencias en un medio lineal}\label{sec:consecuencias}

\begin{proposicion}[Carga ligada en un medio homogéneo]\label{prop:rhoP-rhof}
En el interior de un dieléctrico \LHI{},
\begin{equation}\label{eq:rhoP-rhof}
  \rhoP=-\frac{\chi}{1+\chi}\,\rhof=-\Big(1-\frac{\eo}{\varepsilon}\Big)\rhof .
\end{equation}
En particular, si no hay carga libre en su volumen tampoco hay carga ligada de volumen: la carga de
polarización se concentra en las superficies.
\end{proposicion}

\begin{proof}
Por \eqref{eq:chi}, \eqref{eq:D-epsilon} y \eqref{eq:gauss-D}, con $\chi$ y $\varepsilon$ constantes,
$\rhoP=-\divg\vP=-\eo\chi\,\divg\vE=-\dfrac{\eo\chi}{\varepsilon}\divg\vD=-\dfrac{\chi}{1+\chi}\,\rhof$.
\end{proof}

\begin{proposicion}[Un único medio que llena todo el espacio]\label{prop:un-medio}
Si todo el espacio está ocupado por un único dieléctrico \LHI{} de permitividad $\varepsilon$, entonces
\begin{equation}\label{eq:coulomb-epsilon}
  \divg\vE=\frac{\rhof}{\varepsilon},\qquad \rot\vE=\vcero,\qquad
  \vE(\vr)=\frac{1}{4\pi\varepsilon}\int\frac{\rhof(\vr\,')\,\uv R}{R^2}\,\dd V' :
\end{equation}
el campo es el del vacío con $\eo$ sustituida por $\varepsilon$, y por tanto $\er$ veces menor.
\end{proposicion}

\begin{proof}
Con $\varepsilon$ constante, $\divg\vE=\divg\vD/\varepsilon=\rhof/\varepsilon$. El campo $\varepsilon\vE=\vD$ tiene divergencia $\rhof$ y rotacional
nulo, igual que el campo de Coulomb de la carga libre en el vacío multiplicado por $\eo$; por el
\cref{lem:helmholtz} ambos coinciden.
\end{proof}

\begin{trampa}[¡Solo si hay un único medio!]
Las expresiones $\divg\vE=\rhof/\varepsilon$, $\oiint\vE\cdot\dS=\qfenc/\varepsilon$ y $\vE=\frac{1}{4\pi\varepsilon}\int\rhof\,\uv R/R^2\,\dd V'$ \emph{no son
ciertas en general}. Si $\varepsilon$ varía de un punto a otro,
\[
  \rhof=\divg(\varepsilon\vE)=\varepsilon\,\divg\vE+\vE\cdot\grad\varepsilon ,
\]
y en una interfase, donde $\varepsilon$ cambia bruscamente, $\grad\varepsilon$ es singular: ese término extra es precisamente la
carga de polarización acumulada en la frontera. Por eso en las fronteras se trabaja con $\vD$, cuyas
condiciones de contorno no dependen del material (\cref{sec:frontera}).
\end{trampa}

Los dos ejemplos canónicos son una carga libre en un dieléctrico infinito, que reduce su campo en el factor
$\er$ (\cref{prob:carga-en-dielectrico}), y el condensador plano relleno de dieléctrico, cuya capacidad
aumenta en el mismo factor, $C=\er C_0$ (\cref{prob:condensador-dielectrico}). De ahí procede el método
experimental clásico para medir $\er$: comparar la capacidad de un condensador lleno y vacío.
